\documentclass[11pt,a4paper]{article}

\usepackage[utf8]{inputenc}
\usepackage[T1]{fontenc}

% Deutsche Bezeichner ohne Abhängigkeit von einer Babel-Sprachinstallation.
\renewcommand{\tablename}{Tabelle}
\renewcommand{\figurename}{Abbildung}
\renewcommand{\partname}{Teil}
\renewcommand{\contentsname}{Inhaltsverzeichnis}

% Moderate Seitenränder ohne das geometry-Paket.
\setlength{\textwidth}{15.5cm}
\setlength{\textheight}{23.5cm}
\setlength{\oddsidemargin}{0.5cm}
\setlength{\evensidemargin}{0.5cm}
\setlength{\topmargin}{-0.5cm}

\title{Technische Dokumentation\\\large Chat-Assistant-AI -- Backend}
\author{}
\date{}

\begin{document}

\maketitle

\noindent
Diese Dokumentation gliedert sich in zwei Teile. \textbf{Teil~I} beschreibt die
geplanten Arbeitspakete -- den verbindlichen Funktionsumfang des Systems. \textbf{Teil~II}
dokumentiert begleitende Maßnahmen, die nicht zum geplanten Arbeitspaket-Umfang gehören,
im Zuge der Umsetzung jedoch notwendig wurden: Fehlerbehebungen und Verbesserungen der
zugrunde liegenden Infrastruktur.

\tableofcontents

% ======================================================================================
\part{Arbeitspakete (geplanter Funktionsumfang)}
% ======================================================================================

\section{Session Memory: Persistente Konversationsspeicherung (AP~3A)}

\subsection{Zielsetzung}

Der Konversationsverlauf einer Chat-Sitzung muss dauerhaft und unabhängig von der
Laufzeit des Anwendungsprozesses verfügbar sein. Anforderung ist, dass jeder
Gesprächsbeitrag pro Sitzung verlustfrei persistiert wird, der Verlauf einen
Prozessneustart übersteht und über eine Schnittstelle abrufbar bleibt.

\subsection{Architektur: Zweischichtiges Speichermodell}

Die Konversationsspeicherung ist in zwei klar getrennte Schichten mit
unterschiedlicher Lebensdauer aufgeteilt:

\begin{itemize}
  \item \textbf{Persistente Schicht (SQLite):} Eine relationale Tabelle hält jeden
        Gesprächsbeitrag dauerhaft und bildet die maßgebliche Datenquelle
        (\emph{Single Source of Truth}).
  \item \textbf{Flüchtige Schicht (In-Memory-Puffer):} Ein prozessgebundener
        Ringpuffer hält die jüngsten Beiträge einer Sitzung für den schnellen
        Aufbau des Modell-Prompts vor.
\end{itemize}

Die Entscheidung für eine durable Persistenzschicht unterhalb des In-Memory-Puffers
begründet sich durch mehrere technische Vorteile. Die relationale Ablage gewährleistet
die Dauerhaftigkeit des Verlaufs über Prozessneustarts hinweg und stellt eine eindeutige,
konsistente Datenquelle für alle abhängigen Komponenten bereit. Der vorgelagerte
In-Memory-Puffer vermeidet zugleich, dass für den Aufbau jedes Prompts ein
Datenbankzugriff erforderlich ist, und entkoppelt damit die latenzkritische
Antwortgenerierung von der persistenten Ablage. Beide Schichten ergänzen sich additiv:
Der Puffer fungiert als Lese-Cache über der Datenbank, während die Datenbank die
Verbindlichkeit der Daten sicherstellt.

\subsection{Datenmodell}

Die persistente Schicht wird durch die Tabelle \texttt{chat\_messages} realisiert.
Jeder Datensatz repräsentiert genau einen Gesprächsbeitrag und ist über einen
Fremdschlüssel der zugehörigen Sitzung zugeordnet.

\begin{table}[h]
\centering
\begin{tabular}{lll}
\hline
\textbf{Spalte} & \textbf{Typ} & \textbf{Beschreibung} \\
\hline
\texttt{id}          & INTEGER (PK) & Fortlaufender Primärschlüssel \\
\texttt{session\_id} & TEXT (FK)    & Verweis auf \texttt{chat\_sessions.id}, indiziert \\
\texttt{role}        & TEXT         & Rolle des Beitrags (\texttt{user} / \texttt{assistant}) \\
\texttt{content}     & TEXT         & Inhalt des Beitrags \\
\texttt{created\_at} & TIMESTAMP    & Erstellungszeitpunkt \\
\hline
\end{tabular}
\caption{Schema der Tabelle \texttt{chat\_messages}.}
\end{table}

Der Index auf \texttt{session\_id} ermöglicht das effiziente Auslesen des Verlaufs
einer einzelnen Sitzung. Die monoton steigende \texttt{id} dient zugleich als
verlässliches Sortierkriterium für die chronologische Reihenfolge der Beiträge.

\subsection{Persistenzschicht}

Der Datenbankzugriff ist in einer dedizierten Persistenzkomponente
(\texttt{MessageStore}) gekapselt. Diese Trennung von der Anwendungslogik wurde
gewählt, um die Datenzugriffslogik isoliert testbar zu halten und die Endpunkte von
Datenbankdetails zu entlasten. Die Komponente stellt zwei Operationen bereit:

\begin{itemize}
  \item \texttt{append\_turn} -- persistiert einen vollständigen Austausch
        (Nutzer- und Assistentenbeitrag) in einer einzelnen Transaktion.
  \item \texttt{load\_recent} -- liest die jüngsten Beiträge einer Sitzung,
        begrenzt durch ein Limit und in chronologischer Reihenfolge
        (ältester zuerst), sodass das Ergebnis unmittelbar zum Wiederaufbau des
        Puffers oder Prompts geeignet ist.
\end{itemize}

\begin{verbatim}
class MessageStore:
    @staticmethod
    async def append_turn(db, session_id, user_message, assistant_message):
        db.add(ChatMessage(session_id=session_id, role="user", content=user_message))
        db.add(ChatMessage(session_id=session_id, role="assistant",
                           content=assistant_message))
        await db.commit()

    @staticmethod
    async def load_recent(db, session_id, limit):
        stmt = (
            select(ChatMessage)
            .where(ChatMessage.session_id == session_id)
            .order_by(ChatMessage.id.desc())
            .limit(limit)
        )
        rows = (await db.execute(stmt)).scalars().all()
        return [{"role": r.role, "content": r.content} for r in reversed(rows)]
\end{verbatim}

Die Begrenzung der gelesenen Beiträge über ein Limit hält die Antwortgröße auch bei
langlaufenden Sitzungen beschränkt und vermeidet unbegrenzte Datenbankabfragen.

\subsection{Integration in den Anfrageablauf}

Die Persistenzschicht ist an drei Stellen des Anfrageablaufs eingebunden:

\paragraph{Persistierung der Beiträge.} Beide Chat-Endpunkte -- die
Sofortantwort sowie der tokenweise Stream -- schreiben nach erfolgreicher
Antwortgenerierung den Austausch sowohl in den In-Memory-Puffer als auch in die
Datenbank. Der Schreibvorgang in die Datenbank ist bewusst fehlertolerant
ausgelegt: Da die Antwort zu diesem Zeitpunkt bereits erzeugt wurde, wird ein
transienter Datenbankfehler protokolliert und unterdrückt, statt die Anfrage
oder den laufenden Stream abzubrechen. Diese Auslegung priorisiert die
Verfügbarkeit der Antwort gegenüber der Verbindlichkeit eines einzelnen
Schreibvorgangs.

\paragraph{Wiederherstellung des Puffers (Rehydrierung).} Trifft eine Anfrage auf
eine bestehende Sitzung, deren Puffer im aktuellen Prozess nicht vorliegt -- etwa
nach einem Neustart -- wird der Puffer aus den zuletzt persistierten Beiträgen der
Sitzung rekonstruiert. Dadurch steht der Konversationskontext auch nach einem
Prozesswechsel unmittelbar für den Prompt-Aufbau zur Verfügung. Liegt der Puffer
bereits im Speicher vor, entfällt dieser Schritt, sodass die Datenbank die
maßgebliche Datenquelle bleibt und der Puffer als Cache wirkt.

\paragraph{Abruf des Verlaufs.} Der Endpunkt zum Abruf des Sitzungsverlaufs liest
die Beiträge ausschließlich aus der Datenbank. Damit ist der zurückgelieferte
Verlauf unabhängig vom Zustand des flüchtigen Puffers und auch nach einem
Prozessneustart korrekt.

\subsection{Qualitätssicherung}

Die Funktionalität ist durch automatisierte Tests abgesichert. Diese umfassen die
isolierte Prüfung der Persistenzschicht (Schreib-/Lese-Zyklus, chronologische
Reihenfolge, Begrenzung und Sitzungsbezug der Ergebnisse) sowie Integrationstests
auf Endpunktebene. Letztere verifizieren insbesondere die Persistierung eines
Beitrags und seinen anschließenden Abruf, die Rekonstruktion des Puffers nach einem
simulierten Prozessneustart sowie die Beständigkeit des Verlaufs bei Verlust des
flüchtigen Puffers.

% ======================================================================================
\part{Begleitende Maßnahmen: Bugfixes \& Infrastruktur-Upgrades}
% ======================================================================================

\noindent
Die folgenden Maßnahmen gehören nicht zum geplanten Arbeitspaket-Umfang. Sie wurden im
Zuge der Umsetzung erforderlich -- als Behebung einer Sicherheitslücke beziehungsweise
als Verbesserung der zugrunde liegenden Datenbank-Infrastruktur -- und sind daher
gesondert dokumentiert.

\section{Datenbank-Migrationen mit Alembic (Infrastruktur-Upgrade)}

\subsection{Ausgangslage}

Das Datenbankschema wurde ursprünglich beim Anwendungsstart über die
SQLAlchemy-Funktion \texttt{create\_all()} erzeugt. Dieses Verfahren legt
ausschließlich \emph{fehlende} Tabellen an; bestehende Tabellen werden nie verändert.
Strukturänderungen an einer bereits befüllten Datenbank -- etwa das Hinzufügen einer
Spalte -- sind damit nicht ohne Datenverlust möglich.

\subsection{Entscheidung und Begründung}

Es wurde entschieden, die Schemaverwaltung auf das Migrationswerkzeug \emph{Alembic}
umzustellen. Diese Entscheidung begründet sich durch mehrere strukturelle Vorteile:

\begin{itemize}
  \item \textbf{Versionierung:} Jede Schemaänderung wird als nummeriertes,
        nachvollziehbares Migrationsskript festgehalten.
  \item \textbf{Reversibilität:} Migrationen besitzen eine Vorwärts- und eine
        Rückwärtsrichtung und können kontrolliert zurückgenommen werden.
  \item \textbf{Reproduzierbarkeit:} Aus demselben Migrationsstand entsteht in jeder
        Umgebung dasselbe Schema.
  \item \textbf{Sichere Evolution:} Bestehende, befüllte Datenbanken lassen sich
        strukturell weiterentwickeln, ohne Daten zu verlieren.
\end{itemize}

\subsection{Umsetzung}

Alembic übernimmt die Rolle der maßgeblichen Schema-Autorität. Da die Anwendung eine
asynchrone Datenbank-Engine verwendet, ist die Migrationsumgebung entsprechend
asynchron ausgelegt. Eine Besonderheit von SQLite ist die stark eingeschränkte
Unterstützung von \texttt{ALTER TABLE}: Spalten- oder Constraint-Änderungen sind nicht
direkt möglich. Dies wird über den sogenannten \emph{Batch-Modus} gelöst, in dem die
betroffene Tabelle intern neu aufgebaut und der Datenbestand umkopiert wird.

Die Migrationen sind als geordnete Kette organisiert: Eine Basis-Migration bildet das
vollständige bestehende Schema ab, darauf folgen inkrementelle Migrationen für jede
weitere Änderung (siehe Abschnitt~\ref{sec:ownership}).

\subsection{Integration in Anwendung und Tests}

Beim Start der Anwendung wird das Schema automatisch auf den neuesten Migrationsstand
gebracht. Dieser Vorgang läuft in einem separaten Arbeits-Thread, da die
Migrationsumgebung die asynchrone Engine über einen eigenen Ereignis-Loop steuert.

Die Testsuite ist von diesem Mechanismus bewusst entkoppelt: Tests erzeugen ihr Schema
unmittelbar aus den Modellen auf flüchtigen In-Memory-Datenbanken. Die Migrationen
betreffen somit ausschließlich die persistente Datei-Datenbank, während die Tests
schnell und unabhängig bleiben. Ein zusätzlicher Test stellt sicher, dass die
Migrationskette eine leere Datenbank vollständig bis zum aktuellen Stand aufbaut.

\section{Session-Autorisierung: Ownership-Scoping (Sicherheits-Bugfix)}
\label{sec:ownership}

\subsection{Problembeschreibung}

Das Sitzungsmodell besaß keine Zuordnung zu einem Benutzer. Die Sitzungs-Endpunkte
verlangten zwar eine gültige Authentifizierung, prüften jedoch keine Eigentümerschaft
an der angefragten Sitzung. Dadurch konnte jeder authentifizierte Benutzer den Verlauf
einer beliebigen Sitzung abrufen, sofern ihm deren Bezeichner bekannt war.

Da dieser Bezeichner kein Geheimnis ist -- er erscheint in Antworten und im
clientseitigen Zustand -- handelt es sich um eine Autorisierungslücke: eine fehlende
Zugriffskontrolle auf Objektebene. Das Risiko wurde als hoch eingestuft, da kein
weiterer Schutz zwischen einem gültigen Benutzerkonto und fremden Sitzungsdaten lag.

\subsection{Lösung}

Das Sitzungsmodell wurde um eine Eigentümer-Referenz erweitert (Fremdschlüssel auf den
Benutzer). Beim Anlegen einer Sitzung wird der aufrufende Benutzer als Eigentümer
hinterlegt; bei jedem Zugriff wird die Sitzung auf den jeweiligen Eigentümer
eingeschränkt.

Ein Zugriff auf eine fremde Sitzung wird als \emph{nicht gefunden} (HTTP~404)
beantwortet, nicht als \emph{verboten} (HTTP~403). Diese Wahl ist bewusst getroffen:
Sie verhindert, dass der Endpunkt die Existenz fremder Sitzungen offenlegt. Die
Einschränkung gilt einheitlich für den Abruf des Verlaufs und für das Fortsetzen einer
Sitzung im Chat.

Die zugehörige Schemaänderung -- das Hinzufügen der Eigentümer-Spalte -- wurde über
eine Alembic-Migration eingespielt und ist damit versioniert und reproduzierbar.

\subsection{Qualitätssicherung}

Das Verhalten ist durch Tests abgesichert, die insbesondere den Zugriff durch einen
anderen Benutzer prüfen: Sowohl der Abruf des Verlaufs als auch das Fortsetzen einer
fremden Sitzung müssen mit der Antwort \emph{nicht gefunden} abgewiesen werden, während
der rechtmäßige Eigentümer regulären Zugriff erhält.

\end{document}
